\chapter{Valeurs, complétude et blocs de la hiérarchie spectrale}
\label{chap:torres}

L'opérateur à deux couches, ses moments et le semi-groupe \(\mathfrak S=\langle90,120\rangle\) sont définis une seule fois dans le \cref{sec:jerarquia-espectral-infinita-rev7}. Ce chapitre expose des valeurs, démontre la complétude reconstructive et sépare les blocs globaux. Les mêmes entiers \(90,120\) apparaissent également comme cardinalités dans l'incidence hexade--octade, mais cette coïncidence ne transporte pas l'incidence \(\iota_{6,8}\) dans la carte angulaire et ne transforme pas ses configurations en moments de masse.

\section{Premières évaluations de la hiérarchie}

Les indices
\[
 90,120,180,210,240,270,360,420
\]
constituent une sélection de premières hauteurs de \(\mathfrak S\), puisque
\[
 180=2\cdot90,\quad210=90+120,\quad240=2\cdot120,
\]
\[
 270=3\cdot90,\quad360=3\cdot120,
 \quad420=2(90+120).
\]
Les moments \(s_n=q_-^n-q_+^n\) et l'interface génératrice \(\mathcal K_{\mathrm{ang}}^{90,120}=(s_{90},s_{120})\) sont utilisés ici selon la convention déjà fixée ; ils ne sont pas redéfinis. L'ensemble des valeurs est déterminé par \(A\), \(\Cstar\) et la récurrence quadratique du lecteur.

\begin{table}[htbp]
\centering
\caption{Moments orientés de la hiérarchie globale.}
\label{tab:torres}
\begin{tabular}{@{}rr@{}}
\toprule
\(n\)&\(s_n\)\\
\midrule
90  & \(2.9516943928982796\times10^{-4}\)\\
120 & \(1.9683511250294992\times10^{-5}\)\\
180 & \(8.7345871869508335\times10^{-8}\)\\
210 & \(5.8181313057291163\times10^{-9}\)\\
240 & \(3.8754674969305336\times10^{-10}\)\\
270 & \(2.5814553607509555\times10^{-11}\)\\
360 & \(7.6293257871406076\times10^{-15}\)\\
420 & \(3.3850663628348497\times10^{-17}\)\\
\bottomrule
\end{tabular}
\end{table}

\begin{proposition}[Caractère global des moments]
Les quantités \(s_n\) de la \cref{tab:torres} sont des invariants de la carte angulaire fixée. Une espèce fournit des coordonnées entières qui pondèrent ces invariants, mais ne modifie pas \(q_\pm\) et n'introduit pas de tour propre.
\end{proposition}
\begin{proof}
Chaque \(s_n\) est une fonction de \(q_-\) et de \(q_+\). Le théorème d'inversion des canaux démontre que l'application conjointe \((A,\Cstar)\mapsto(q_+,q_-)\) est injective. Aucune étiquette d'espèce n'apparaît dans ces définitions.
\end{proof}

\begin{theorem}[Complétude reconstructive de la hiérarchie]
\label{thm:completitud-reconstructiva-torres} Soit \(n_1<n_2<\cdots\) l'énumération croissante sans répétitions de \(\mathfrak S\). Pour
\[
 a_k:=s_{n_k}=q_-^{n_k}-q_+^{n_k}>0,
\]
tout rapport positif \(\rho\in\RR_{>0}\) admet un développement absolument convergent dans la hiérarchie :
\[
 \boxed{\rho=
 \exp\!\left(\sum_{k\geq1}\nu_k a_k\right)}
 \qquad\text{pour une certaine suite }(\nu_k)_{k\geq1}\subset\ZZ.
\]
\end{theorem}

\begin{proof}
Puisque \(0<q_+<q_-<1\), on a
\[
 0<a_k<q_-^{n_k}.
\]
Les \(n_k\) sont des multiples de \(30\) supérieurs ou égaux à \(90\) ; donc,
\[
 \sum_{k\geq1}a_k
 \leq\sum_{m\geq3}q_-^{30m}
 =\frac{q_-^{90}}{1-q_-^{30}}<\infty,
 \qquad a_k\longrightarrow0.
\]
Fixons \(r_0=\log\rho\) et une règle de départage pour l'entier le plus proche. Définissons récursivement
\[
 \nu_k=\operatorname{nint}\!\left(\frac{r_{k-1}}{a_k}\right),
 \qquad
 r_k=r_{k-1}-\nu_k a_k.
\]
La propriété de l'entier le plus proche donne
\[
 |r_k|\leq\frac{a_k}{2}\longrightarrow0.
\]
L'identité télescopique
\[
 r_0=\sum_{j=1}^{N}\nu_j a_j+r_N
\]
implique \(r_0=\sum_{j\geq1}\nu_j a_j\). La convergence est absolue, puisque
\[
 \begin{aligned}
 \sum_{k\geq1}|\nu_k|a_k
 &=\sum_{k\geq1}|r_{k-1}-r_k|\\
 &\leq |r_0|+\frac{a_1}{2}
 +\frac12\sum_{k\geq2}(a_{k-1}+a_k)<\infty.
 \end{aligned}
\]
L'exponentiation donne l'énoncé.
\end{proof}

\begin{remark}[Portée et contrôle négatif]
Le théorème établit une surjectivité reconstructive du langage de tours sur \(\RR_{>0}\) : il démontre la résolution du codomaine, et non la sélection physique. L'algorithme reçoit \(r_0=\log\rho\) ; si \(\rho\) est une grandeur cible, ses coefficients constituent donc une reconstruction dépendant de cette cible. La provenance prospective exige une autre flèche,
\[
 \text{espèce}\dashrightarrow\text{extension}
 \longrightarrow\text{route}\longrightarrow\text{registre}
 \longrightarrow\text{signature},
\]
et ne se déduit pas de ce développement. La conclusion serait invalidée si les \(a_k\) ne tendaient pas vers zéro, si la borne du résidu cessait d'être satisfaite ou si la série des erreurs n'était pas sommable.
\end{remark}

\section{Blocs globaux de l'action de masse}

L'action de masse utilisera le triplet
\[
 \Dfour,\qquad s_{120},\qquad \zCorr=135\Dfour^2.
\]
Leurs valeurs sont calculées une seule fois. Les coordonnées \(k,\nu_{120},\nu_{270}\) d'une signature précisent le nombre d'interventions de chaque bloc dans l'exposant ; elles n'en modifient pas la valeur.

Il convient de distinguer deux objets qui partagent l'indice \(270\) :
\[
 \boxed{\zCorr=135\Dfour^2
 =1.66922699663643\times10^{-6}}
\]
et
\[
 \boxed{\sSpec=q_-^{270}-q_+^{270}
 =2.58145536075096\times10^{-11}.}
\]
Le premier est un invariant quadratique construit à partir de \(\Dfour\) ; le second est un moment spectral du semi-groupe. La coïncidence de l'indice n'autorise pas leur identification.

\section{Combinaison d'ordre supérieur}

La combinaison globale
\[
 \delta\Dfour^{\rm tower}
 =\frac{s_{180}}{24}-\frac{s_{240}}{54}
 -24s_{360}+\frac76s_{420}
\]
a pour valeur
\[
 \delta\Dfour^{\rm tower}
 =3.632051471908759123889070967\times10^{-9}.
\]
La coordonnée de comparaison utilisée pour sélectionner les coefficients est définie comme la donnée scalaire
\[
 \delta\Dfour
 :=3.632051471908972784661641720\times10^{-9}.
\]
Par conséquent, le résidu de la combinaison par rapport à cette coordonnée est
\[
 \delta\Dfour^{\rm tower}-\delta\Dfour
 =-2.1366077257075\times10^{-22}.
\]
La petite quantité est donc une différence entre deux objets, et non la valeur de la combinaison. Tous les termes appartiennent à la même hiérarchie et leurs coefficients sont globaux. La coordonnée de comparaison étant intervenue dans la sélection des coefficients, cette égalité est reconstructive. La combinaison n'est pas ajoutée une nouvelle fois à la formule électronique qui emploie déjà \(\Dfour\), car cette substitution compterait deux fois la même correction.
